% Clase 18 --- Las fuerzas sobre la espira y el balance de energía (§5.1).
% "Van a tener un tramo de fuerza acá, un tramo de fuerza acá, un tramo de
%  fuerza acá... F2, fíjense que intenta oponerse al movimiento, intenta frenar."
% Misma geometría que la figura anterior: acá sólo se agregan las fuerzas.
% Sentidos calculados: con I horaria, en el lado izquierdo (dentro del campo) la
% corriente va hacia ARRIBA, y F = I D y x (-B z) = -I D B x, o sea hacia la
% IZQUIERDA: frena. Los tramos horizontales dan fuerzas verticales opuestas.
\begin{center}
\begin{tikzpicture}[scale=1]

  % región con campo entrante
  \fill[figblue!8] (-3.9,-1.9) rectangle (0,1.9);
  \draw[figgray, line width=0.7pt] (-3.9,-1.9) rectangle (0,1.9);
  \foreach \xx in {-3.5,-2.9,-2.3,-1.7,-1.1,-0.5} {
    \foreach \yy in {-1.5,1.5} {
      \node[figblue, scale=0.6] at (\xx,\yy) {$\otimes$};
    }
  }
  \node[figlblaux, anchor=south west] at (-3.9,1.95) {$\vec B$ entrante};

  % la espira y la corriente inducida
  \draw[figline, line width=1.5pt] (-1.3,-1.05) rectangle (1.5,1.05);
  \draw[figvecres, line width=1pt] (-0.5,1.05) -- (0.6,1.05);
  \draw[figvecres, line width=1pt] (1.5,0.4) -- (1.5,-0.4);
  \draw[figvecres, line width=1pt] (0.6,-1.05) -- (-0.5,-1.05);
  \draw[figvecres, line width=1pt] (-1.3,-0.4) -- (-1.3,0.4);
  \node[figlbl, figred, anchor=west] at (1.65,1.05) {$I = \dfrac{BDv}{R}$};

  % fuerzas: las de los tramos horizontales se cancelan
  \draw[figvec, line width=1.2pt] (-0.7,1.05) -- (-0.7,1.75);
  \node[figlbl, figblue, anchor=south] at (-0.7,1.78) {$\vec F_1$};
  \draw[figvec, line width=1.2pt] (-0.7,-1.05) -- (-0.7,-1.75);
  \node[figlbl, figblue, anchor=north] at (-0.7,-1.78) {$\vec F_3$};

  % la fuerza del lado interior: frena
  \draw[figvecres, line width=1.6pt] (-1.3,0) -- (-2.5,0);
  \node[figlbl, figred, anchor=south] at (-2.5,0.12) {$\vec F_2 = I D B$};

  % velocidad
  \draw[figvecres, line width=1.3pt] (1.7,0) -- (2.7,0);
  \node[figlbl, figred, anchor=south] at (2.2,0.06) {$\vec v$};
  \node[figlblaux, anchor=west, align=left] at (1.75,-0.95)
    {alguien tiene que tirar\\[-0.25em] para mantener $v$};

  \node[figlbl, anchor=north, align=center] at (-0.6,-2.35)
    {$P_{\text{dis}} = R I^2 = \dfrac{B^2D^2v^2}{R}
     \;=\; F_2\,v = P_{\text{ext}}$};

\end{tikzpicture}\\[0.5em]
{\small\itshape Una vez que hay corriente, el mismo campo actúa sobre ella:
$\vec F = I\vec L\times\vec B$ en cada tramo. Los dos tramos horizontales que
están dentro del campo llevan corriente en sentidos opuestos y dan fuerzas
verticales que se cancelan (además de tender a deformar la espira, no a
moverla). Todo el efecto neto queda en el lado interior, y su fuerza
$F_2 = I D B = B^2D^2v/R$ apunta \textbf{hacia adentro del campo}: se opone al
movimiento. Eso no es casualidad ni una coincidencia de signos, es Lenz otra
vez, y el balance energético lo confirma: la espira disipa $RI^2$ por efecto
Joule, así que para mantener $v$ constante alguien tiene que entregar
exactamente esa potencia, y en efecto $F_{\text{ext}}v = B^2D^2v^2/R = RI^2$. Si
nadie tira, la energía cinética se va convirtiendo en calor y el sistema
\emph{se frena solo}: es el principio del \textbf{freno magnético}, que no tiene
contacto ---no se gastan pastillas--- pero tampoco puede detener del todo, ya
que al anularse $v$ desaparece la fuerza.}
\end{center}
